\documentclass{article}
\usepackage[T1]{fontenc}
\usepackage{lmodern}
\usepackage{graphicx}
\usepackage{amsmath,amssymb}
\usepackage[a4paper,margin=2cm]{geometry}
\usepackage[absolute,overlay]{textpos}
\setlength{\TPHorizModule}{1mm}
\setlength{\TPVertModule}{1mm}
\usepackage[indonesian]{babel}
\usepackage{hyperref}
\setlength{\emergencystretch}{2em}
% unit: Halaman 2

\begin{document}

% === Halaman 2 (python/native) ===

2

Materi Matematika untuk Kriptorafi:

1.

Teori Bilangan


- Integer dan aritmetika modulo


- Algoritma euclidean


- Kekongruenan


- Relatif prima


- Balikan (invers) modulo


- Bilangan prima

2.

Probabilitas dan Statistik

3.

Kompleksitas algoritma

4.

Teori Informasi

5.

Aljabar abstrak 

No. 1 s/d 3 sudah dipelajari di dalam 

kuliah Matematika Diskrit dan Probabilitas 

dan Statistik

No 5 akan dibahas pada materi ECC 

(Elliptic Curve Cryptography)

\end{document}
